\begin{helper}
Let $\mathcal F\in H(k,k,D)$, let $D\ge 2$, let $g$ be an edge of
$\mathcal F$, let $v\in g$, and let $x_1$ be one of the $k-1$ fresh
vertices in the split-edge construction of
\noderef{KUniformSplitEdgeHypergraph}.  Then
$\operatorname{SplitEdge}_k(\mathcal F,g,v,x_1)$ is also a member of
$H(k,k,D)$.
\end{helper}
\begin{proof}
By \noderef{HypergraphClass}, it is enough to prove $k$-uniformity,
$k$-simplicity, and maximum degree at most $D$ for the split hypergraph.
The $k$-uniformity statement is exactly \noderef{KUniformSplitEdgeUniformity}.

For $k$-simplicity, use only this uniformity.  If $a$ and $b$ are two
distinct edge labels in the split hypergraph, then
$a\cap b\subseteq a$.  Since $a$ has exactly $k$ vertices by the already
proved uniformity, $|a\cap b|\le |a|=k$.  This is the required
\noderef{TSimpleHypergraph} condition with $t=k$.

It remains to check degrees, using \noderef{HypergraphDegree} and
\noderef{MaxDegreeAtMost}.  For an old vertex $w\ne v$, the split
operation preserves exactly the old incident labels: old labels other than
$g$ are unchanged, the modified copy of $g$ still contains $w$ precisely
when $g$ originally contained $w$, and the one new edge contains only the
old vertex $v$ together with fresh vertices.  Hence the degree of $w$ is
unchanged and is at most $D$.

For the old vertex $v$, the old label $g$ no longer contains $v$, while the
one new edge does contain $v$; every other old label contains $v$ exactly
when it did before.  Thus the degree of $v$ is also unchanged, and is at
most $D$ by the original maximum-degree hypothesis.

Finally let $y$ be a fresh vertex.  The only split-edge labels that can
contain $y$ are the modified old label $g$, when $y=x_1$, and the single new
edge.  Therefore $y$ has degree at most $2$, which is at most $D$ by
hypothesis.  All vertices have degree at most $D$, so the split hypergraph
lies in $H(k,k,D)$.
\end{proof}
